\documentclass[11pt]{article}
\usepackage{amsmath,amssymb,bm}
\usepackage[margin=2.4cm]{geometry}
\usepackage{booktabs}

\title{The RAYEN head on the BiLSTM gap imputer\\
\large size-aware balance split, as implemented}
\author{}
\date{}

\begin{document}
\maketitle

\section{What it is}

The imputer is a bidirectional LSTM whose output passes through a fixed,
non-learned, differentiable map before the loss is taken. The map is a RAYEN
head, generalised from a single timestep to the whole gap trajectory.

\paragraph{Pipeline.}
\begin{align*}
\mathbf{D} &= \mathrm{Linear}_{256\to 6}\bigl(\mathrm{BiLSTM}(x,\ \mathrm{mask})\bigr)
 && \text{network output, standardised units}\\
\mathbf{F} &= (\mathbf{I}+\mathbf{D})\odot\bm{\sigma}_y+\bm{\mu}_y
 && \text{raw fill in MW; generally infeasible}\\
\mathbf{P} &= \mathrm{RayenHead}(\mathbf{F})
 && \text{feasible output}\\
\mathcal{L} &= \mathrm{MSE}(\mathbf{P},\mathbf{Y})
 && \text{loss on the \emph{mapped} output}
\end{align*}

Three facts about the head:
\begin{itemize}
\item It has \textbf{no learnable parameters}. It is a fixed map.
\item In the code it is \textbf{not} inside \texttt{BiLSTMImputer}.
      \texttt{model.py} returns the raw $\mathbf{D}$; the caller applies the map
      (\texttt{plan1\_train.py:66} in training, \texttt{plan1\_bench.py:66} at eval).
\item It acts on the \textbf{whole $36\times6$ trajectory at once}, not per
      timestep. One scalar step length governs all $216$ numbers. This is the
      difference from a single-step RAYEN head, and it is what makes the head
      behave conservatively in practice (Section \ref{sec:measured}).
\end{itemize}

\paragraph{The idea in one sentence.}
The head starts from a trajectory that is guaranteed feasible, walks toward the
one the network asked for, and stops at the first physical limit.

\section{Notation}

\begin{center}
\begin{tabular}{llp{7.6cm}}
\toprule
symbol & type & meaning \\
\midrule
\multicolumn{3}{l}{\emph{sizes and indices}}\\
$N$ & $=36$ & unknown steps in the gap (3 h at 5 min) \\
$K$ & $=6$ & dispatch channels \\
$t$ & $1..N$ & step index inside the gap \\
$i$ & $1..K$ & channel index \\
$j$ & $1..N{+}1$ & index over consecutive differences (includes both seams) \\
\midrule
\multicolumn{3}{l}{\emph{known inputs to the head}}\\
$\mathbf{p}^{L},\mathbf{p}^{R}$ & $\mathbb{R}^{K}$ & true dispatch one step before / after the gap (MW), held pinned \\
$nd_t$ & $\mathbb{R}$ & balance target at step $t$ (MW); see note below \\
$\mathbf{c}$ & $\mathbb{R}^{K}$ & per-channel cap (MW) \\
$\bm{\rho}^{+},\bm{\rho}^{-}$ & $\mathbb{R}^{K}_{\ge0}$ & up / down ramp limit per 5-min step (MW), stored positive \\
$\mathbf{s}$ & $\mathbb{R}^{K}$ & balance signs $(1,1,1,1,-1,1)^{\!\top}$; charging is a load. Note $s_i^2=1$ \\
\midrule
\multicolumn{3}{l}{\emph{trajectories, all $\mathbb{R}^{N\times K}$ in MW}}\\
$\mathbf{I}$ & & straight interpolation from $\mathbf{p}^{L}$ to $\mathbf{p}^{R}$ \\
$\mathbf{F}$ & & raw fill: what the network asked for. Infeasible in general \\
$\mathbf{A}$ & & \textbf{anchor}: the guaranteed-feasible fallback \\
$\mathbf{r}$ & & \textbf{direction}: $\mathbf{F}-\mathbf{A}$, flattened so it preserves balance \\
$\mathbf{P}$ & & \textbf{the head's output}. Satisfies every constraint, so $\mathbf{s}^{\!\top}\mathbf{P}_t=nd_t$ exactly \\
$\mathbf{X}$ & & the \emph{iterate} inside the Step-1 repair loop. Starts at $\mathbf{I}$, ends at $\mathbf{A}$. Infeasible in between. Not the output \\
$\mathbf{w}$ & & size-aware weights, $w_{t,i}=|\cdot|+1$ \\
\midrule
\multicolumn{3}{l}{\emph{step length}}\\
$\alpha$ & $\mathbb{R}$ & position along the ray. $\alpha=0$ gives $\mathbf{A}$; $\alpha=1$ gives $\mathbf{A}+\mathbf{r}$, the network's request \emph{after} balance correction. Note $\mathbf{A}+\mathbf{r}\neq\mathbf{F}$ unless $\mathbf{F}$ already balanced \\
$\alpha^{\star}$ & $[0,1]$ & \textbf{the chosen step length}: how far the head actually walks. Set by box and ramp only, never by balance \\
$\alpha_{\mathrm{hi}},\alpha_{\mathrm{lo}}$ & $\mathbb{R}$ & largest $\alpha$ before hitting the upper / lower box wall \\
$\alpha_{\mathrm{up}},\alpha_{\mathrm{dn}}$ & $\mathbb{R}$ & largest $\alpha$ before hitting the up / down ramp wall \\
$\alpha_{\mathrm{soc}}$ & $\mathbb{R}$ & same for the battery-energy wall (available, off by default) \\
\midrule
\multicolumn{3}{l}{\emph{other}}\\
$\delta_t$ & $\mathbb{R}$ & balance residual of the \emph{iterate}, $nd_t-\mathbf{s}^{\!\top}\mathbf{X}_t$. Nonzero during Step 1, zero once converged \\
$M$ & $\mathbb{N}$ & anchor refinement cycles: $40$ in training, $160$ at eval \\
$\bm{\mu}_y,\bm{\sigma}_y$ & $\mathbb{R}^{K}$ & target scaler mean / scale \\
$\odot$ & & elementwise product \\
$\eta,\Delta t,C_{\mathrm{eff}}$ & & $\sqrt{0.834}$, $1/12$ h, $4535.75$ MWh: battery constants \\
\bottomrule
\end{tabular}
\end{center}

\paragraph{Which $nd$.} Two different quantities are called net demand.
The \emph{network input} is demand-side,
$\text{demand}-\text{wind}-\text{solar}-\text{curtailments}$ (feature column 6).
The \emph{head's balance target} $nd_t$ is supply-side,
$nd_t=\mathbf{s}^{\!\top}\mathbf{P}^{\text{true}}_t$
(\texttt{plan1\_train.supply\_nd}; \texttt{gw.truth\_mw @ SIGN} at eval). They
differ by the interconnector, mean $\pm500$ MW. The supply-side value is used
because it is attainable by six fuels; the demand-side one is not.

\section{The head, in four steps}

\subsection*{Step 1 --- the fallback $\mathbf{A}$}

Take the straight line between the pinned boundaries,
$\mathbf{I}_t=\mathbf{p}^{L}+\tfrac{t}{N+1}(\mathbf{p}^{R}-\mathbf{p}^{L})$, and
repair it into feasibility by cycling three projections $M$ times:
\begin{equation}
\mathbf{X}\ \leftarrow\ \Pi_{\mathrm{box}}\bigl(\Pi_{\mathrm{ramp}}(\Pi_{\mathrm{bal}}(\mathbf{X}))\bigr),
\qquad \mathbf{X}^{(0)}=\mathbf{I},
\qquad \mathbf{A}=\mathbf{X}^{(M)} .
\end{equation}
$\mathbf{X}$ is the \emph{iterate}, and it is infeasible until the loop settles.
This matters for reading Section \ref{sec:sizeaware}: the balance operator's
argument is $\mathbf{X}$, not the head's output $\mathbf{P}$, so its residual
$\delta_t$ is genuinely nonzero. It is what the operator exists to remove.
\begin{itemize}
\item $\Pi_{\mathrm{bal}}$: snap the signed sum onto $nd_t$ (Section \ref{sec:sizeaware}).
\item $\Pi_{\mathrm{ramp}}$: one forward clamp from $\mathbf{p}^{L}$, one backward clamp from $\mathbf{p}^{R}$.
\item $\Pi_{\mathrm{box}}$: clip to $[0,c_i]$.
\item $\Pi_{\mathrm{soc}}$: scale both battery channels to fit the pack. \textbf{Off} on this path.
\end{itemize}
$\mathbf{A}$ is built from $\mathbf{p}^{L},\mathbf{p}^{R},nd$ only. It does not
see the network and carries no gradient.

\paragraph{Why $M$ is large.} $\Pi_{\mathrm{bal}}$ drives $\delta$ to zero in one
shot, but $\Pi_{\mathrm{ramp}}$ and $\Pi_{\mathrm{box}}$ then move the numbers
again and break balance, so the next cycle starts from a fresh nonzero
$\delta$. The three operators settle against each other slowly. On the hardest
of 60 test windows, $\delta=\max_t|nd_t-\mathbf{s}^{\!\top}\mathbf{X}_t|$ at the
end of each cycle runs:
\begin{center}
\begin{tabular}{lrrrrrrr}
\toprule
cycles $M$ & 0 & 1 & 5 & 20 & 40 & 80 & 160 \\
\midrule
$\delta$ (MW) & $602.0$ & $76.0$ & $61.9$ & $30.0$ & $12.1$ & $2.1$ & $0.061$ \\
\bottomrule
\end{tabular}
\end{center}
Training uses $M=40$ and eval uses $M=160$, which is where the $12.1$ MW
train-time imbalance in Section \ref{sec:measured} comes from.

\subsection*{Step 2 --- the direction $\mathbf{r}$}

Walking from $\mathbf{A}$ straight at $\mathbf{F}$ would break balance, because
$\mathbf{F}$ does not satisfy it. So first flatten the direction:
\begin{equation}
\mathbf{r}=\mathbf{F}-\mathbf{A},
\qquad
\mathbf{r}_t\ \leftarrow\ \mathbf{r}_t-\frac{\mathbf{s}^{\!\top}\mathbf{r}_t}{\sum_i w_{t,i}}\,(\mathbf{w}_t\odot\mathbf{s}),
\qquad w_{t,i}=|A_{t,i}|+1 .
\label{eq:tangent}
\end{equation}
After this, $\mathbf{s}^{\!\top}\mathbf{r}_t=0$ for every $t$: moving along
$\mathbf{r}$ changes the mix but not the signed sum. The check is one line, using $s_i^2=1$:
\[
\mathbf{s}^{\!\top}\mathbf{r}_t-\frac{\mathbf{s}^{\!\top}\mathbf{r}_t}{\sum_i w_{t,i}}\sum_i w_{t,i}s_i^{2}=0 .
\]
The weights come from $\mathbf{A}$, not $\mathbf{F}$, so \eqref{eq:tangent} is a
fixed linear operator and the gradient to the network is clean.

\subsection*{Step 3 --- the step length $\alpha^{\star}$}

\paragraph{What $\alpha^{\star}$ is not for.} Balance is already finished. Step 2
made $\mathbf{s}^{\!\top}\mathbf{r}_t=0$, so
$\mathbf{s}^{\!\top}(\mathbf{A}+\alpha\mathbf{r})_t=nd_t$ at \emph{every}
$\alpha$, including $\alpha=1$. In particular $\mathbf{A}+\mathbf{r}$ meets
demand exactly (measured: worst residual $9.2\times10^{-4}$ MW over 40 test
windows). $\alpha^{\star}$ exists only because $\mathbf{A}+\mathbf{r}$ can still
break the \emph{other} constraints: measured on the same windows it violates the
box in $39/40$ windows by up to $1541$ MW and the ramp in $33/40$ by up to
$2613$ MW. Step 3 is about box and ramp, nothing else.

Along $\mathbf{y}(\alpha)=\mathbf{A}+\alpha\mathbf{r}$ each of those constraints
is a straight line in $\alpha$, so each gives a simple ratio
\emph{(slack at the anchor) / (rate of approach)}:
\begin{align}
\text{box:}\quad
&\alpha_{\mathrm{hi}}=\min_{r_{t,i}>0}\frac{c_i-A_{t,i}}{r_{t,i}},
&&\alpha_{\mathrm{lo}}=\min_{r_{t,i}<0}\frac{0-A_{t,i}}{r_{t,i}},
\label{eq:wall-box}\\[3pt]
\text{ramp:}\quad
&\alpha_{\mathrm{up}}=\min_{\Delta r_{j,i}>0}\frac{\rho^{+}_i-\Delta A_{j,i}}{\Delta r_{j,i}},
&&\alpha_{\mathrm{dn}}=\min_{\Delta r_{j,i}<0}\frac{-\rho^{-}_i-\Delta A_{j,i}}{\Delta r_{j,i}},
\label{eq:wall-ramp}
\end{align}
where $\Delta\mathbf{A}=\mathrm{diff}([\mathbf{p}^{L};\mathbf{A};\mathbf{p}^{R}])$
and $\Delta\mathbf{r}=\mathrm{diff}([\mathbf{0};\mathbf{r};\mathbf{0}])$ are the
$N{+}1$ consecutive differences. Padding $\mathbf{r}$ with zeros says the head
does not move the pinned boundaries, which is what puts both seams under the
same ramp limit as the interior. Then
\begin{equation}
\boxed{\ \alpha^{\star}=\operatorname{clip}\bigl(\min(\alpha_{\mathrm{hi}},\alpha_{\mathrm{lo}},\alpha_{\mathrm{up}},\alpha_{\mathrm{dn}}),\ 0,\ 1\bigr)\ }
\end{equation}
The clip at $1$ stops the head overshooting past $\mathbf{F}$. The clip at $0$
matters more than it looks; see Section \ref{sec:measured}.

\subsection*{Step 4 --- the output}
\begin{equation}
\boxed{\ \mathbf{P}=\mathbf{A}+\alpha^{\star}\mathbf{r}\ }
\end{equation}

\section{What ``size-aware'' changes}
\label{sec:sizeaware}

Only one thing: how a balance correction is shared out across the six channels.

\emph{Read $\mathbf{X}$ below as the mid-loop iterate, not the head's output.}
Its residual $\delta_t=nd_t-\mathbf{s}^{\!\top}\mathbf{X}_t$ is nonzero (up to
$602$ MW at $M=0$); the output $\mathbf{P}$ has $\delta_t=0$ by construction.
The balance snap moves along $\mathbf{w}_t\odot\mathbf{s}$:
\begin{equation}
\Pi_{\mathrm{bal}}(\mathbf{X})_t=\mathbf{X}_t+\frac{\delta_t}{\sum_i w_{t,i}}\,(\mathbf{w}_t\odot\mathbf{s}),
\label{eq:bal}
\end{equation}
which lands exactly on $nd_t$ for \emph{any} positive $\mathbf{w}$, by the same
$s_i^2=1$ identity as before. The arms differ only in $\mathbf{w}$:
\begin{center}
\begin{tabular}{lll}
\toprule
arm & weights & effect \\
\midrule
\texttt{baseline}, \texttt{cost} & $w_{t,i}=1$ & every channel moves the same MW \\
\texttt{size\_aware} & $w_{t,i}=|X_{t,i}|+1$ & each channel moves in proportion to its size \\
\bottomrule
\end{tabular}
\end{center}

\paragraph{Why it matters.} \texttt{gas\_steam} typically runs at single-digit
MW. Under the equal split, a $600$ MW residual hands it $100$ MW, many times the
channel's own output. Under the size-aware split, \texttt{coal\_brown}
(thousands of MW) absorbs the bulk and \texttt{gas\_steam} moves a proportionate
sliver. The $+1$ MW floor keeps $\mathbf{w}>0$ so \eqref{eq:bal} stays exact when
a channel sits at zero.

Geometrically, \eqref{eq:bal} is the oblique projection onto the balance
hyperplane along $\mathbf{w}\odot\mathbf{s}$. Only $\mathbf{w}\equiv1$ is the
orthogonal, minimum-norm projection; size-aware deliberately gives that up to
keep the fuel mix ratios.

\paragraph{Where the flag acts.} Two places, both inside the head: the anchor's
balance snap in Step 1, and the tangent removal \eqref{eq:tangent} in Step 2.
Nothing else in the model changes.

\section{What is guaranteed}

\begin{description}
\item[Balance --- exact, at any $\alpha$.] $\mathbf{s}^{\!\top}\mathbf{r}_t=0$,
      so $\mathbf{s}^{\!\top}\mathbf{P}_t=\mathbf{s}^{\!\top}\mathbf{A}_t$. The
      output inherits the anchor's balance; the only error is the anchor's own
      residual, which depends on $M$.
\item[Box and ramp --- exact.] Each is a straight line in $\alpha$ and
      $\alpha^{\star}$ is the smallest distance to any of them, so every
      constraint holds on the whole segment $[0,\alpha^{\star}]$.
\item[Differentiable.] $\mathbf{A}$ is constant, $\mathbf{r}$ is affine in
      $\mathbf{F}$, $\alpha^{\star}$ is a min of smooth ratios. Gradients reach
      the network through both $\mathbf{r}$ and $\alpha^{\star}$.
\item[SOC --- not enforced.] The $\alpha_{\mathrm{soc}}$ wall exists in the code
      but \texttt{soc=False} is the default and neither the training nor the
      benchmark script overrides it. Measured slack on every test window anyway.
\item[Conservative by design.] One scalar $\alpha^{\star}$ for all $N\times K$
      cells. A single binding cell anywhere caps the entire trajectory.
\end{description}

\subsection*{The bias introduced by a single $\alpha^{\star}$}

Sharing one step length across all $216$ cells is not neutral. It biases the
head in three distinct ways.

\paragraph{1. The head is invariant to the network's gain.} When a box wall
binds, $\alpha_{\mathrm{lo}}=A_{t,i}/(-r_{t,i})$ at the binding cell, so
\[
\alpha^{\star}\,r_{t,i}=-A_{t,i}
\qquad\text{independently of }\|\mathbf{r}\| .
\]
If the network doubles every output, $\mathbf{r}$ doubles, $\alpha^{\star}$
halves, and the applied move $\alpha^{\star}\mathbf{r}$ is unchanged. The total
movement is set by the anchor's slack at the tightest cell, not by anything the
network asks for. The network can change the \emph{shape} of its deviation but
not the \emph{amount}, and cannot learn to push harder.

\paragraph{2. One channel's limit shrinks all six.} The binding cell is
\texttt{battery\_discharging} in $30/40$ windows for \texttt{size\_aware}. That
battery hitting its zero floor scales down the \texttt{coal\_brown},
\texttt{hydro} and gas predictions by the identical factor. Nothing physical
couples them; it is an artifact of the shared scalar.

\paragraph{3. The gradient is throttled with the output.}
$\partial\mathbf{P}/\partial\mathbf{F}$ carries a factor of order
$\alpha^{\star}$, so the network trains at roughly $6\%$ signal. Visible in the
histories: \texttt{baseline} ($\alpha^{\star}\!\approx\!0$) has a validation MSE
spanning $0.041$ over 17 epochs, essentially flat because the network barely
reaches the output, against $0.194$ for \texttt{size\_aware}.

\paragraph{Counterpoint.} At present the throttle helps. Mean MAE against truth
over 40 windows: anchor $\mathbf{A}$ $50.7$ MW, head output $\mathbf{P}$
$56.4$ MW, but $\mathbf{A}+\mathbf{r}$ (the full request, $\alpha=1$)
$491.4$ MW for \texttt{baseline} and $316.3$ MW for \texttt{size\_aware}. The
network's balance-corrected request is far worse than the fallback, so a small
$\alpha^{\star}$ is currently suppressing a bad network rather than a good one.
Removing the throttle will make the numbers worse before it makes them better.

\section{As implemented, and what it measures}
\label{sec:measured}

\paragraph{Configuration.} Anchor cycles $M=40$ in training,
$M=160$ at eval; \texttt{soc=False} both; balance target supply-side; float32 in
training, float64 at eval. The $M$ mismatch is real: at $M=40$ the output carries
up to $12.1$ MW of balance residual, at $M=160$ it is $0.061$ MW. Reconstruction
MAE is unaffected ($53.60$ vs $53.61$).

\paragraph{Feasibility of the trained \texttt{size\_aware} output}
(60 test 3-h gaps, $M=160$): worst balance residual $6.1\times10^{-2}$ MW, worst
ramp overshoot $0$, most negative dispatch $1.8\times10^{-15}$ MW, SOC-infeasible
windows $0/60$. The construction holds.

\paragraph{$\alpha^{\star}$ collapses.} Measured on 40 test gaps with the trained
checkpoints:
\begin{center}
\begin{tabular}{lrrr}
\toprule
 & \texttt{baseline} & \texttt{cost} & \texttt{size\_aware} \\
\midrule
mean $\alpha^{\star}$          & $0.017$ & $0.006$ & $0.097$ \\
median $\alpha^{\star}$        & $0$     & $0$     & $0.056$ \\
windows with $\alpha^{\star}=0$ & $32/40$ & $38/40$ & $3/40$ \\
binding wall                   & box-lo $39$ & --- & box-lo $34$, ramp-up $6$ \\
\bottomrule
\end{tabular}
\end{center}

\paragraph{Mechanism.} The anchor's last inequality step is the box clip, so
$\mathbf{A}$ routinely sits \emph{on} the lower face, $A_{t,i}=0$ exactly (about
$26$ of $216$ cells per window). If any such cell has $r_{t,i}<0$, then
\eqref{eq:wall-box} gives
\begin{equation}
\alpha_{\mathrm{lo}}\le\frac{0-A_{t,i}}{r_{t,i}}=0
\quad\Longrightarrow\quad \alpha^{\star}=0
\quad\Longrightarrow\quad \mathbf{P}=\mathbf{A}.
\end{equation}
One cell out of $216$ zeroes the single global step length, and the network's
entire contribution is dropped. For \texttt{baseline} and \texttt{cost} this
happens in most windows, so their output \emph{is} the projected interpolation.

\paragraph{Why size-aware helps.} Its proportional split stops the anchor from
pushing mid-sized channels down onto the zero face: \texttt{hydro} goes
$10.8\%\to0\%$ of cells at zero, \texttt{battery\_charging} $27.7\%\to0\%$. Fewer
zero cells, fewer triggers, so $\alpha^{\star}=0$ drops from $32/40$ to $3/40$.
It does not remove the conservatism: \texttt{gas\_steam} and \texttt{gas\_ocgt}
genuinely sit at zero $44\%$ and $25\%$ of the time, and the median surviving
fraction of the network's deviation is still $5.6\%$, i.e.\ $445$ MW of model
deviation admitted as $36$ MW of movement.

\section{Code map}

\begin{center}
\begin{tabular}{ll}
\toprule
what & where \\
\midrule
caps, ramps, $\eta$, pack size & \texttt{ml/check\_caps.py} \\
$\Pi_{\mathrm{bal}}$ (the size-aware split) & \texttt{constraints.py:104} \texttt{\_balance\_project} \\
$\Pi_{\mathrm{ramp}}$ & \texttt{constraints.py:123} \texttt{\_ramp\_project} \\
$\Pi_{\mathrm{soc}}$ & \texttt{constraints.py:141} \texttt{\_soc\_project} \\
anchor loop (Step 1) & \texttt{constraints.py:157} \texttt{cyclic\_project} \\
Steps 2--4 (the head) & \texttt{constraint\_layers.py:60} \texttt{rayen\_traj\_project} \\
network + loss & \texttt{model.py} \\
head applied, training & \texttt{plan1\_train.py:66} \texttt{forward\_mapped} \\
head applied, eval & \texttt{plan1\_bench.py:66} \texttt{score} \\
readable single-file version & \texttt{sizeaware\_constraints\_reference.py} \\
\bottomrule
\end{tabular}
\end{center}

\noindent
\emph{Two implementation notes.} \texttt{ramp\_tube}, \texttt{bridge\_feasible}
and \texttt{\_ramp\_pocs} in \texttt{constraints.py} are analysis helpers and are
never called on this path. And \texttt{constraints.py} caches
$\mathbf{c},\bm{\rho}^{\pm}$ as float32 tensors, so even the float64 path runs
with limits off by up to $1.95\times10^{-4}$ MW; harmless, but it is why a clean
float64 reimplementation disagrees at that level.

\end{document}
